\begin{abstract}
 Network routers that enforce Quality-of-Service (QoS) guarantees must decide, at every clock cycle, which expiring packet of information to transmit, even when the value of the packet is \textit{unknown} until it is processed. We frame this problem as the \textit{Online Packet Scheduling with Deadlines} (OPSD) problem under Partial Feedback: packets arrive at every clock cycle, with different \textit{deadlines}, but the weights are only observed after execution. Under a stochastic assumption on the unknown weights, we explore different variants of the OPSD problem with bandit feedback. We establish a connection between our setting and the \textit{sleeping bandits} problem, and set our learning goal to $\alpha$-regret minimization. We provide algorithms with provable $\alpha$-regret guarantees under different spans of slackness, distinguishing systems allowing for randomization and systems that do not. In every scenario, our algorithms achieve an $\alpha$-regret upper bound of $\widetilde{\mathcal{O}}\left(\sqrt{KT}\right)$, matching the lower bound for the standard bandit setting. In the practically relevant case of $2$-bounded deadline instances, where the deadline is set at most one clock cycle away from the arrival, our deterministic algorithm achieves the provably tightest possible competitive ratio. Remarkably, when the number of distinct packet types $K\ge 2$ is finite, it is possible to break the well-established $\Phi = \frac{1+\sqrt{5}}{2}$ competitive ratio barrier and attain a tighter competitive ratio $\theta_K$ ranging in $[\sqrt{2}, \Phi)$.
\end{abstract}


% \begin{abstract}
% Network routers enforcing Quality-of-Service must transmit packets with hard deadlines, yet a packet’s value may only be revealed after transmission. We study \emph{Online Packet Scheduling with Deadlines} under \emph{bandit feedback}: at each time step, the scheduler selects one feasible packet and observes only that packet’s realised weight. Assuming packets belong to $K$ observable types with unknown stochastic distributions, we introduce the \emph{K-OPSD} problem and evaluate performance via the \emph{$\alpha$-regret}, which combines competitive analysis with statistical regret. We relate K-OPSD to \emph{sleeping bandits} and design algorithms across slackness regimes, for both deterministic and randomized systems. In all settings, we obtain $\widetilde{\mathcal{O}}\left(\sqrt{KT}\right)$ $\alpha$-regret, matching standard stochastic bandit lower bounds. In the practically relevant \emph{2-bounded} case, we give a deterministic algorithm with a tight competitive ratio. Remarkably, when the number of types $K\ge 2$ is finite, we show that it is possible to break the classical $\Phi=\frac{1+\sqrt{5}}{2}$ competitive ratio barrier, achieving a tighter $\theta_K \in [\sqrt{2}, \Phi)$.
% \end{abstract}



  % Buffer management in Quality of Service (QoS) switches is a key problem in modern network routing optimization. QoS devices, such as network routers, may face a large amount of data traffic and only have a limited amount of time to select which portion of that should be executed, prioritizing the most important pieces of information. Existing literature framed this problem in the \textit{Online Packet Scheduling with Deadlines} framework: an algorithm is faced with repeated arrivals of unit-cost packets, and each packet has its own \textit{deadline} and \textit{weight}. Without foreseeing the future, the algorithm is then asked to choose one packet per time unit and gains its weight. In this work, we initiate the study of a new variant of this problem, which accounts for realistic scenarios with partial information. In particular, we assume that weights can only be observed once the packet is processed, a common scenario when extremely low latency does not allow for a proper weight estimation in real-time. However, \textit{metadata} (such as the packet's origin or information typology) are usually available. By considering the weights stochastic and iid distributed when conditioned on the metadata, we show how learning can be performed to obtain competitive ratio plus regret type of guarantees.
  %Network routers that enforce Quality-of-Service (QoS) guarantees must decide, at every clock cycle, which expiring packet of information to transmit, even when its value is \textit{unknown} until it is processed. We frame this problem as Online Packet Scheduling with Deadlines under Partial Feedback: packets arrive at every clock cycle, with different \textit{deadlines} and \textit{metadata} (such as their origin or information type), but the actual weight is only revealed at service time. Under the assumption that weights are stochastic and iid distributed when conditioned on the metadata, we explore different variants of the traditional Online Packet Scheduling with Deadlines problem under the lens of learning with bandit feedback. We show an important connection between our setting and the \textit{Sleeping Bandits Problem}, and set our learning goal to $\alpha$-regret minimization. We provide algorithms with provable $\alpha$-no-regret guarantees under different spans of slackness, with a major focus on the practically relevant $2$-bounded regimes, where the deadline is set at most one clock cycle away from the arrival. As customary in existing literature, we explicitly distinguish between systems allowing for randomization and systems that do not, the latter being the most common case in real-world applications. In every scenario, our algorithms achieve an $\alpha$-regret upper bound in the order of $\sqrt{T}$, which is the known lower bound for standard bandit setting. In the special case of $2$-bounded instances, our deterministic algorithm also achieves the provably tightest possible competitive ratio: remarkably, when the number of distinct metadata $K\ge 2$ is finite, it is possible to break the well-established $\Phi$ barrier for the competitive ratio in Online Packet Scheduling with Deadlines, and attain a $\theta_K$ competitive ration spanning in $[\sqrt{2}, \Phi)$.

% I guess there should be more emphasis on our contributions in the abstract: 

% - lower bound (maybe saying "We provide algorithms with provable near-optimal $\alpha$-no-regret guarantees ..."?, also, what are regimes of slackness?  the problem is that we have "near-optimal" guarantees only for deterministic algorithms in the 2-bounded regime, and that is also our only lower bound. Maybe
%  then i guess, after you say "relevant 2-bounded regimes." We say something like, in this setting, we prove matching upper and lower bound on the alpha-regret, using a UCB adaptation of the EDF algorithm
% I'm trying to add something but let's not make it too long


% - upper bounds: what is the conclusion? i.e. learning is possible since it is a sublinear regret, matching the lower bound! added some stuff

% - algorithmic ideas: new algorithm that solves the open problem of theta_k cr, what is the main idea/intuition of the algo that makes it work? UCB-adaptation of potential-based algorithms in the randomised setting, etc 

%This does not belong to the abstract in my opinion

%it's too long, we have to prune something
